% !PRINTED 47
\weberpage{47}
%<<CONT
and setting $x = 0, 1, 1:\varkappa^2$ in this identical equation, it
follows that:
\weq{16}{\begin{aligned}
\sqrt{x_1x_2x_3} &= -\dfrac{c}{b} \\
\sqrt{(1-x_1)(1-x_2)(1-x_3)} &= \dfrac{b+a}{b} \\
\sqrt{(1-\varkappa^2x_1)(1-\varkappa^2x_2)(1-\varkappa^2x_3)} &= \dfrac{b+\varkappa^2a}{b}.
\end{aligned}}

For the determination of the signs in these equations we obtain
further from (13), if we replace $x_1$ by $x_2$ and by $x_3$
and multiply the results:
\[ x_1x_2x_3(a+bx_1)(a+bx_2)(a+bx_3) = -c^3\sqrt{f(x_1)}\sqrt{f(x_2)}\sqrt{f(x_3)} \]
and (15) gives for $x = -a:b$
\[ (a+bx_1)(a+bx_2)(a+bx_3) = \dfrac{c^2}{b}(b+a)(b+\varkappa^2a). \]
If one uses in addition the square of the first equation (16),
there follows from this
\[ \sqrt{f(x_1)}\sqrt{f(x_2)}\sqrt{f(x_3)} = -\dfrac{c}{b^3}(b+a)(b+\varkappa^2a). \]

But the same result is given by the multiplication of the three
equations (16), and the signs in these formulae are therefore so
chosen that the product of the left-hand sides gives the value of
$\sqrt{f(x_1)}$ required by the differential equation (1). A further
determination of the signs is, from the nature of the case, not
possible in (16), and these equations serve for the unique
definition of $\sqrt{x_1}, \sqrt{1-x_1}, \sqrt{1-\varkappa^2x_1}$, when the
$\sqrt{x_2}, \sqrt{1-x_2}, \sqrt{1-\varkappa^2x_2}, \sqrt{x_3}, \sqrt{1-x_3}, \sqrt{1-\varkappa^2x_3}$ are
presupposed as given.

In order to bring the equations (16) into the customary form, we
first make the numerator of $b$ in (14) rational, and obtain
\[ b = -\dfrac{(x_2-x_3)(1-\varkappa^2x_2x_3)}{\sqrt{x_2}\sqrt{(1-x_3)(1-\varkappa^2x_3)} + \sqrt{x_3}\sqrt{(1-x_2)(1-\varkappa^2x_2)}}, \]
further, by (14):
\[ \begin{aligned}
b+a &= -\sqrt{(1-x_2)(1-x_3)}\left\{\sqrt{x_2(1-x_2)}\sqrt{1-\varkappa^2x_3} - \sqrt{x_3(1-x_3)}\sqrt{1-\varkappa^2x_2}\right\} \\
b+\varkappa^2a &= -\sqrt{(1-\varkappa^2x_2)(1-\varkappa^2x_3)}\left\{\sqrt{x_2(1-\varkappa^2x_2)}\sqrt{1-x_3} - \sqrt{x_3(1-\varkappa^2x_3)}\sqrt{1-x_2}\right\},
\end{aligned} \]
%CONT>>
